\documentclass{article} % say
\usepackage{amsmath}
\usepackage{gensymb}
\usepackage{hyperref}
\usepackage{tikz}
\usepackage{tkz-euclide}
\usetikzlibrary {babel,angles,quotes,positioning}
\usepackage[utf8]{inputenc}
\usepackage[russian]{babel}
\begin{document}

\textbf{Задача}: Найти все корни уравнения
\begin{equation}
  x \sqrt{2x^2 - 511} - \sqrt{x} \sqrt[4]{2x^2 - 511} = 12
\end{equation}

\clearpage

\textbf{Решение}: Заменим $a = \sqrt{x} \sqrt[4]{2x^2 - 511}$, тогда:

\begin{equation}
  a^2 - a = 12
\end{equation}

Откуда:

\begin{equation}
  \left[
    \begin{aligned}
      \sqrt{x} \sqrt[4]{2x^2 - 511} & = -3 \\
      \sqrt{x} \sqrt[4]{2x^2 - 511} & = 4
    \end{aligned}
  \right.
  \label{eq:union}
\end{equation}

В первом случае решений нет. Докажем это. Возведём обе части в четвёртую степень и сделаем замену $t = x^2$.

\begin{equation}
  t (2t - 1) = 81
\end{equation}

Получим корни:

\begin{equation}
  t = \dfrac{511 \pm \sqrt{261769}}{4}
\end{equation}

Имеем два корня: $t_1$ отрицательное число близкое к нулю, $t_2$~--- положительное число, причём $2t_2 - 511 > 0$.

Рассмотрим $t_1$. В этом случае $x = \pm ai$, где $i$~-- мнимая единица, $a>0$. В любом из этих случаев $\sqrt{x} = \mp \sqrt{a}$, т.е. действительное число, а $x^2 - 511 < 0$, а значит $\sqrt[4]{x^2 - 511}$ не может быть действительным. Их произведение не может быть действительным.

Рассмотрим $t_2$. 

Во втором случае уравнения \ref{eq:union}:

\begin{equation}
  \sqrt{x} \sqrt[4]{2x^2 - 511} = 4
\end{equation}

Возведём обе части в четвёртую степень и сделаем замену $t = x^2$

\begin{equation}
\begin{aligned}
  t (2t - 511) & = 256 \\
  t & = \dfrac{511 \pm 513}{4} 
\end{aligned}
\end{equation}

Отсюда: $x = \{16, -\frac{i}{4}\}$.

\textbf{Источник}: Репетиционное тестирование ЦТ 2023-1, вариант 3, задание B17 \url{https://tinyurl.com/54jakpdm} 

\end{document}


